\chapter{The \prg Runtime Interface}
\label{app:api}

This appendix is the working reference for the \prg framework calls used
throughout the book, together with the calling convention, the cartridge tool, and
the constants you will reach for most often. It catalogues the public interface
declared in \fname{components/prg32/include/prg32.h} and documented in the
platform's ABI notes \citep{prg32repo,prg32abi}. A cartridge can call exactly the
functions listed in the portable cartridge ABI (\fname{prg32/abi/prg32\_abi.json};
version 1.6, 139 entries, at the time of writing). A few functions declared in
the header are \emph{resident-only}: they belong to the firmware, and a cartridge
that calls one fails to link with an undefined-symbol error. It is organised by task rather
than alphabetically, so that you can find a call by what you are trying to do.

\section{The calling convention}

Every framework function follows the standard \riscv calling convention, which is
what allows assembly games and C games to call the same functions
\citep{prg32abi}. Arguments are passed in \reg{a0}--\reg{a7} in order; the return
value comes back in \reg{a0}. Temporaries \reg{t0}--\reg{t6} may be destroyed by a
call; saved registers \reg{s0}--\reg{s11} are preserved across one. The return
address is in \reg{ra} and must be saved before a nested call; the stack pointer
\reg{sp} stays 16-byte aligned.

\begin{center}
\begin{tabularx}{\linewidth}{@{}llX@{}}
\toprule
Register & Role & Preserved across a call? \\
\midrule
\reg{a0}--\reg{a7} & arguments / return value & no \\
\reg{t0}--\reg{t6} & temporaries & no \\
\reg{s0}--\reg{s11} & saved values & yes \\
\reg{ra} & return address & no (save it yourself) \\
\reg{sp} & stack pointer & yes (16-byte aligned) \\
\bottomrule
\end{tabularx}
\captionof{table}{The \riscv register convention as \prg uses it \citep{prg32abi}.}
\end{center}

\section{Lifecycle and timing}

\begin{center}
\begin{tabularx}{\linewidth}{@{}lX@{}}
\toprule
Function & Purpose \\
\midrule
\code{prg32\_ticks\_ms()} & milliseconds since start, for timing and animation \\
\code{prg32\_random\_number(min,max)} & uniform random value in an inclusive range \\
\code{prg32\_gfx\_present()} & display the frame you have drawn \\
\code{prg32\_init()} & initialise the runtime (resident-only; embedded firmware builds) \\
\bottomrule
\end{tabularx}
\captionof{table}{Lifecycle and timing calls.}
\end{center}

\section{Input}

All buttons arrive as a single bitmask. Player two occupies the same layout shifted
into higher bits (\code{PRG32\_P2\_*}) \citep{prg32repo}.

\begin{center}
\begin{tabularx}{\linewidth}{@{}lX@{}}
\toprule
Function & Purpose \\
\midrule
\code{prg32\_input\_read()} & full player 1 / player 2 bitmask \\
\code{prg32\_input\_read\_player(p)} & one player's bits, normalised to the low bits \\
\code{prg32\_input\_read\_menu()} & both joysticks merged, for menu navigation \\
\code{prg32\_input\_wait\_released(mask)} & block until the selected bits are released \\
\bottomrule
\end{tabularx}
\captionof{table}{Input calls.}
\end{center}

\begin{center}
\begin{tabular}{@{}lr l lr@{}}
\toprule
Constant & Bit & & Constant & Bit \\
\midrule
\code{PRG32\_BTN\_LEFT}  & \code{0x01} & & \code{PRG32\_BTN\_A}      & \code{0x10} \\
\code{PRG32\_BTN\_RIGHT} & \code{0x02} & & \code{PRG32\_BTN\_B}      & \code{0x20} \\
\code{PRG32\_BTN\_UP}    & \code{0x04} & & \code{PRG32\_BTN\_SELECT} & \code{0x40} \\
\code{PRG32\_BTN\_DOWN}  & \code{0x08} & & (\code{PRG32\_BTN\_START} alias) & \code{0x40} \\
\bottomrule
\end{tabular}
\captionof{table}{Player-one button bits \citep{prg32repo}.}
\end{center}

\section{Graphics}

The viewport is 320 by 200 within a 320-by-240 display; colours are RGB565
\citep{prg32repo}.

\begin{center}
\begin{tabularx}{\linewidth}{@{}lX@{}}
\toprule
Function & Purpose \\
\midrule
\code{prg32\_gfx\_clear(color)} & fill the viewport with one colour \\
\code{prg32\_gfx\_pixel(x,y,color)} & set one pixel \\
\code{prg32\_gfx\_rect(x,y,w,h,color)} & filled rectangle \\
\code{prg32\_gfx\_text8(x,y,s,fg,bg)} & draw a text string \\
\code{prg32\_gfx\_present()} & present the finished frame \\
\code{prg32\_gfx\_set\_fullscreen(on)} & draw into the full 320-by-240 display instead of the viewport \\
\bottomrule
\end{tabularx}
\captionof{table}{Core graphics calls.}
\end{center}

\begin{center}
\begin{tabularx}{\linewidth}{@{}lX@{}}
\toprule
Function & Purpose \\
\midrule
\code{prg32\_palette\_set(index,rgb565)} & assign a colour to one of 256 palette entries \\
\code{prg32\_palette\_get(index)} & read a palette entry \\
\code{prg32\_gfx\_clear\_indexed(index)} & fill the viewport with a palette index \\
\code{prg32\_gfx\_rect\_indexed(x,y,w,h,index)} & filled rectangle in a palette index \\
\code{prg32\_gfx\_pixel\_indexed(x,y,index)} & one pixel in a palette index \\
\bottomrule
\end{tabularx}
\captionof{table}{Indexed (palette) graphics calls.}
\end{center}

\begin{center}
\begin{tabularx}{\linewidth}{@{}lX@{}}
\toprule
Function & Purpose \\
\midrule
\code{prg32\_band\_set\_mode(band,mode)} & what a status band shows: nothing, FPS, Wi-Fi, game, debug, custom \\
\code{prg32\_band\_set\_text(band,text)} & custom text for a band \\
\code{prg32\_band\_set\_game\_info(text)} & the game-information line (score, level) \\
\code{prg32\_band\_set\_colors(band,fg,bg)} & band colours \\
\bottomrule
\end{tabularx}
\captionof{table}{The two 20-pixel status bands above and below the viewport
(\code{PRG32\_BAND\_TOP}, \code{PRG32\_BAND\_BOTTOM}).}
\end{center}

\begin{center}
\begin{tabular}{@{}lr l lr@{}}
\toprule
Colour & RGB565 & & Colour & RGB565 \\
\midrule
\code{PRG32\_COLOR\_BLACK} & \code{0x0000} & & \code{PRG32\_COLOR\_YELLOW}  & \code{0xffe0} \\
\code{PRG32\_COLOR\_WHITE} & \code{0xffff} & & \code{PRG32\_COLOR\_CYAN}    & \code{0x07ff} \\
\code{PRG32\_COLOR\_RED}   & \code{0xf800} & & \code{PRG32\_COLOR\_MAGENTA} & \code{0xf81f} \\
\code{PRG32\_COLOR\_GREEN} & \code{0x07e0} & & \code{PRG32\_COLOR\_BLUE}    & \code{0x001f} \\
\bottomrule
\end{tabular}
\captionof{table}{Named colour constants \citep{prg32repo}.}
\end{center}

\section{Sprites and collision}

\begin{center}
\begin{tabularx}{\linewidth}{@{}lX@{}}
\toprule
Function & Purpose \\
\midrule
\code{prg32\_sprite\_hitbox(ax,ay,aw,ah,bx,by,bw,bh)} & non-zero if two rectangles overlap \\
\code{prg32\_sprite\_anim\_frame(now,count,ms)} & which animation frame to show now \\
\code{prg32\_sprite\_draw\_frame(x,y,w,h,frames,frame,transparent)} & draw one frame of an animation \\
\code{prg32\_sprite\_draw\_8x8(x,y,bits,fg,bg)} & draw an 8-by-8 bit sprite \\
\code{prg32\_sprite\_draw\_16x16(x,y,rgb565)} & draw a 16-by-16 RGB565 sprite (also \code{\_24x24}) \\
\code{prg32\_sprite\_anim\_init / \_update / \_draw} & an animated sprite object that advances its own frames \\
\code{prg32\_sprite\_draw\_indexed(x,y,\&spr,frame)} & draw a palette-indexed sprite (1, 2, 4, or 8 bits per pixel) \\
\bottomrule
\end{tabularx}
\captionof{table}{Sprite and collision calls \citep{prg32repo}.}
\end{center}

\section{Playfields and tiles}

The framework provides two scrolling tile layers for larger worlds
\citep{prg32repo}.

\begin{center}
\begin{tabularx}{\linewidth}{@{}lX@{}}
\toprule
Function & Purpose \\
\midrule
\code{prg32\_playfield\_put(layer,tx,ty,id)} & place a tile in the grid \\
\code{prg32\_playfield\_get(layer,tx,ty)} & read a tile from the grid \\
\code{prg32\_playfield\_scroll(layer,x,y)} & scroll a layer to a position \\
\code{prg32\_playfield\_parallax(layer,x\_q8,y\_q8)} & scroll with fixed-point parallax \\
\code{prg32\_playfield\_camera(x,y)} & set the shared camera position \\
\code{prg32\_playfield\_draw\_dual()} & draw both layers \\
\bottomrule
\end{tabularx}
\captionof{table}{Playfield and tile calls.}
\end{center}

\section{Platform-game physics}

The actor struct \code{prg32\_platform\_actor\_t} and its helpers implement
gravity, movement, and tile collision \citep{prg32repo}.

\begin{center}
\begin{tabularx}{\linewidth}{@{}lX@{}}
\toprule
Function & Purpose \\
\midrule
\code{prg32\_platform\_tile\_flags(id,flags)} & give a tile behaviour (solid, hazard, \dots) \\
\code{prg32\_platform\_actor\_init(\&a,layer,x,y,w,h)} & set up an actor \\
\code{prg32\_platform\_actor\_step(\&a,input,move,jump,grav,maxfall)} & one frame of physics; returns state flags \\
\code{prg32\_platform\_camera\_follow(\&a,dzx,dzy)} & scroll to keep the actor in view \\
\bottomrule
\end{tabularx}
\captionof{table}{Platform physics calls. Pass the actor by address (\code{\&}).}
\end{center}

\begin{center}
\begin{tabular}{@{}ll l ll@{}}
\toprule
Tile flag & Meaning & & Actor state & Meaning \\
\midrule
\code{...\_SOLID}    & blocks movement & & \code{...\_ON\_GROUND} & standing on ground \\
\code{...\_PLATFORM} & one-way platform & & \code{...\_HIT\_HEAD}  & bumped a ceiling \\
\code{...\_HAZARD}   & harmful tile & & \code{...\_HAZARD}    & touched a hazard \\
\code{...\_COLLECT}  & collectable & & \code{...\_COLLECT}   & collected an item \\
\bottomrule
\end{tabular}
\captionof{table}{Tile flags (\code{PRG32\_TILE\_FLAG\_*}) and actor state bits
(\code{PRG32\_PLATFORM\_*}) \citep{prg32repo}.}
\end{center}

\section{Console}

\begin{center}
\begin{tabularx}{\linewidth}{@{}lX@{}}
\toprule
Function & Purpose \\
\midrule
\code{prg32\_console\_clear()} & clear the 40-by-25 text console \\
\code{prg32\_console\_putc(ch)} & write one character \\
\code{prg32\_console\_write(s)} & write a string to the console \\
\code{prg32\_console\_hex32(v)} & write a 32-bit value in hexadecimal \\
\bottomrule
\end{tabularx}
\captionof{table}{Console calls.}
\end{center}

\section{Sound}

The mixer has eight voices at 22\,050\,Hz. It drives the I2S amplifier on the
board and host audio on QEMU and PRG32-QT. Samples, instruments, and tracker
music come from the optional audio block of the cartridge \citep{prg32cartfmt}.

\begin{center}
\begin{tabularx}{\linewidth}{@{}lX@{}}
\toprule
Function & Purpose \\
\midrule
\code{prg32\_audio\_note(ch,instr,note,vol,ms)} & play a MIDI note for a duration; returns at once \\
\code{prg32\_audio\_note\_on(ch,instr,note,vol)} / \code{\_note\_off(ch)} & start and stop a sustained note \\
\code{prg32\_audio\_notes(ch,instr,vol,notes,count)} & play a sequence of notes (blocking) \\
\code{prg32\_audio\_play\_sample(id,vol,pitch)} & play a sample from the audio block \\
\code{prg32\_audio\_play\_track(id)} / \code{\_stop\_track()} & start and stop tracker music \\
\code{prg32\_audio\_set\_master\_volume(v)} & overall volume \\
\code{prg32\_audio\_set\_channel\_pan(ch,pan)} & stereo position of a channel \\
\code{prg32\_buzzer\_tone(hz,ms,duty)} & legacy: PWM tone on a passive buzzer, if one is fitted \\
\bottomrule
\end{tabularx}
\captionof{table}{Sound calls. Prefer the \code{prg32\_audio\_*} family; the
buzzer helpers remain for older examples. Consult \fname{prg32\_audio.h} for exact
argument types.}
\end{center}

\section{Text entry}

\begin{center}
\begin{tabularx}{\linewidth}{@{}lX@{}}
\toprule
Function & Purpose \\
\midrule
\code{prg32\_text\_input(buffer,capacity,title)} & run a complete on-screen keyboard dialogue \\
\code{prg32\_keyboard\_init / \_update / \_draw} & the same keyboard, driven from your own loop \\
\bottomrule
\end{tabularx}
\captionof{table}{On-screen keyboard calls.}
\end{center}

\section{Scores}

\begin{center}
\begin{tabularx}{\linewidth}{@{}lX@{}}
\toprule
Function & Purpose \\
\midrule
\code{prg32\_score\_submit(game,player,score)} & record a score under a given name \\
\code{prg32\_score\_submit\_current\_player(game,score)} & record a score under the name saved in setup \\
\code{prg32\_score\_player\_prompt()} & ask the player for a name with the on-screen keyboard \\
\code{prg32\_score\_count(game)} / \code{prg32\_score\_get(game,i,\&s)} & read the local leaderboard \\
\code{prg32\_scoreboard\_show(game,title)} & display the leaderboard \\
\code{prg32\_score\_sync\_remote()} & exchange scores with the configured Cartridge Store \\
\bottomrule
\end{tabularx}
\captionof{table}{Score calls \citep{prg32abi}. Remote scores are described in
Chapter~\ref{ch:store}.}
\end{center}

\section{Multiplayer}

A cartridge joins a room identified by a \emph{signature} string; only
cartridges with the same signature see each other. On the board the relay is
reached over Wi-Fi through the Cartridge Store; QEMU provides an offline stub so
that the same code builds and runs.

\begin{center}
\begin{tabularx}{\linewidth}{@{}lX@{}}
\toprule
Function & Purpose \\
\midrule
\code{prg32\_multiplayer\_available()} & non-zero when the service can be used \\
\code{prg32\_multiplayer\_join(signature,flags)} / \code{\_leave()} & enter and leave a room \\
\code{prg32\_multiplayer\_set\_local\_state(\dots)} & publish this player's position and state \\
\code{prg32\_multiplayer\_tick()} & exchange state; call once per frame \\
\code{prg32\_multiplayer\_get\_peer\_count()} / \code{\_get\_peer(i,\dots)} & read the other players \\
\bottomrule
\end{tabularx}
\captionof{table}{Multiplayer calls; see \fname{prg32\_multiplayer.h} for the
peer structure.}
\end{center}

\section{Cartridge slots and network state}

\begin{center}
\begin{tabularx}{\linewidth}{@{}lX@{}}
\toprule
Function & Purpose \\
\midrule
\code{prg32\_cart\_stored\_count()} & how many of the four slots hold a cartridge \\
\code{prg32\_cart\_get\_slot\_info(slot,\&info)} & name, sizes, and flags of a stored cartridge \\
\code{prg32\_cart\_default\_slot()} & saved default cartridge slot, or $-1$ \\
\code{prg32\_cart\_select\_slot(slot)} / \code{\_select\_default()} & switch to another stored cartridge \\
\code{prg32\_wifi\_current\_mode()} / \code{\_current\_ip()} & Wi-Fi mode and address \\
\bottomrule
\end{tabularx}
\captionof{table}{Selected cartridge-slot and network calls \citep{prg32abi}.}
\end{center}

\section{The command-line tool}

Everything is one Python module, run from the \prg repository root with ESP-IDF
active \citep{prg32repo}. Add \code{--help} to any command for its options.

\begin{center}
\begin{tabularx}{\linewidth}{@{}lX@{}}
\toprule
Command & Purpose \\
\midrule
\code{doctor} & check the toolchain and partition table \\
\code{esp32c6 build | flash | build-and-flash} & resident firmware for the board \\
\code{esp32c6 upload | run | upload-and-run} & send a cartridge to a board (or PRG32-QT) over HTTP \\
\code{esp32c6 screenshot | memory | performance} & capture the screen; analyse memory; fetch test results \\
\code{qemu build | run | build-and-run} & emulator firmware and virtual screen \\
\code{qemu upload} & stage a cartridge into the emulator flash image \\
\code{cartridge build} & build a portable \fname{.prg32} from \fname{.S} or \fname{.c} \\
\code{cartridge summary} & print header, ABI, sizes, features, metadata \\
\code{store attach-metadata | inspect-metadata} & write and read the metadata trailer \\
\code{store discover | list | download} & find a Store and fetch from it \\
\code{store pack-bundle | publish-bundle} & create and submit a Store bundle \\
\code{wifi set} & preconfigure the board's Wi-Fi \\
\code{debug enable | pause | step | resume | speed | state} & control the PRG32-QT debugger \\
\code{abi gen | check} & regenerate and verify the ABI headers (framework developers) \\
\bottomrule
\end{tabularx}
\captionof{table}{\code{python3 -m prg32 \dots}: the unified tool.}
\label{tab:cli}
\end{center}

Useful options of \code{cartridge build} are \code{--portable} (always),
\code{--entry-prefix}, \code{--name}, \code{--out}, \code{--architecture}
(\code{esp32c6} or \code{qemu}), \code{--cart-ram-kib}, \code{--audio-block},
\code{--multiplayer}, and \code{--required-feature}.

\section{The HTTP runtime API}

When the board's firmware is running and reachable, it exposes a small HTTP
interface useful for tooling and labs \citep{prg32repo}. PRG32-QT serves the
same interface and adds the debugger route.

\begin{center}
\begin{tabularx}{\linewidth}{@{}lX@{}}
\toprule
Endpoint & Returns \\
\midrule
\code{GET /api/runtime} & firmware version, cartridge state and RAM size, frame count, last input \\
\code{GET /api/games} & the installed cartridges \\
\code{POST /api/games?slot=cart0} & install a cartridge into a slot \\
\code{GET /api/screenshot.bmp} & the current 320-by-240 framebuffer as a BMP \\
\code{GET /api/performance.json} & the latest performance-test result \\
\code{GET}, \code{POST /api/debug} & PRG32-QT only: debugger state and commands \\
\bottomrule
\end{tabularx}
\captionof{table}{Selected HTTP runtime endpoints \citep{prg32repo,prg32qt}.}
\end{center}
